\chapter*{学生实验}\markboth{学生实验}{学生实验}
\stepcounter{chapter}
\ctexset{section/name={实验,}}
\addcontentsline{toc}{chapter}{学生实验}
\renewcommand{\thefigure}{\arabic{figure}}
\renewcommand{\thetable}{\arabic{table}}

\section{铜和它的化合物的性质}
\subsection*{实验目的}
加深对铜和它的化合物的性质的认识。

\subsection*{实验用品}
试管、试管架、试管夹、滴管、量筒、酒精灯、砂纸。

铜丝、铜片、浓硫酸、稀硫酸、浓硝酸、稀硝酸、稀盐酸、$0.2M$ \ce{AgNO3} 溶液、\ce{CuO} 粉末、\ce{CuSO4.5H2O} 晶体、$0.2M$ \ce{CuSO4} 溶液、$2M$ \ce{NaOH} 溶液、$2M$ 氨水、\ce{BaCl2} 溶液。

\subsection*{实验步骤}
\begin{enumerate}[1.]
  \item 铜的性质
  \begin{enumerate}[(1)]
    \item 取一段铜丝，用砂纸把它的一端磨光，在酒精灯火焰上加热几秒钟后取下，冷却。观察铜丝表面发生了什么变化。写出反应的化学方程式。
    \item 在四个试管里分别加人 \qty{2}{ml} 稀硫酸、浓硫酸、稀硝酸和浓硝酸。再各加入一小块铜片，记下观察到的现象\footnote{做这个实验时，要注意实验室的通风。}。对没有反应发生的试管加微热。再记下观察到的现象。解释这些现象并写出有关的化学方程式。
    \item 在一个试管里加入 \qty{2}{ml} $0.2M$ \ce{AgNO3} 溶液。把一根洁净的铜丝弯成螺旋状，浸入 \ce{AgNO3} 溶液中。观察发生的现象，并判断铜丝上和试管里生成的固体是什么物质。试用实验证明你的判断。
  \end{enumerate}
  \item 铜的化合物的性质
  \begin{enumerate}
    \item 在三个装有少量 \ce{CuO} 粉末的试管里，分别加入 \qty{2}{ml} 稀盐酸、稀硫酸和稀硝酸。观察各个试管里发生的现象；然后分别加热，再观察发生的现象。
    \item 在一个干燥的试管里加入少量 \ce{CuSO4.5H2O} 晶体，把试管口微向下倾，加热硫酸铜晶体，发生什么现象？试管冷却后，加入少量蒸馏水，得到蓝色溶液。向蓝色溶液里滴加 \ce{BaCl2} 溶液，产生白色沉淀。然后加稀硝酸，沉淀不溶解。试说明这一系列的实验现象，并写出反应的化学方程式。
    \item 向盛有 \qty{2}{ml} $0.2M$ \ce{CuSO4} 溶液的试管里滴加 \qty{2}{ml} $2M$ 的 \ce{NaOH} 溶液,观察沉淀的生成。把生成的沉淀分放到三个试管里，向其中的一个试管滴加稀硫酸，向另一个试管滴加 \ce{NaOH} 溶液,给第三个试管加热。观察并解释这些现象。
    \item 在试管里加入 \qty{2}{ml} $0.2M$ \ce{CuSO4} 溶液，滴加 \qty{2}{ml} $2M$ \ce{NaOH} 溶液，再滴加浓氨水，观察蓝色沉淀的生成和溶解。
    
    在生成的溶液中加入约 \qty{2}{ml} $2M$ \ce{NaOH} 溶液，观察有没有反应发生。
  \end{enumerate}

  解释上述现象，并写出反应的离子方程式。
\end{enumerate}

\begin{Theorem}{讨论}
  \begin{enumerate}[1.]
    \item 铜能跟浓硫酸起反应，为什么不能跟稀硫酸起反应？如果用铁代替铜跟浓、稀硫酸起反应，会发生什么现象？
    \item 实验 1. (3) 说明了什么？如果用 \ce{Hg(NO3)2} 或 \ce{Pb(NO3)2} 代替 \ce{AgNO3}，会发生什么现象？
  \end{enumerate}
\end{Theorem}

\section{实验习题}
\subsection*{实验目的}
巩固有关铁、铜的性质和络合物的知识。
\subsection*{实验习题}
\begin{enumerate}[1.]
  \item 现有一瓶溶液，怎样通过实验来确定它是氯化铁？
  \item 用实验区别 \ce{Fe^3+} 和 \ce{[Fe(CN)6]^3-} 络离子。
  \item 怎样从溶液里分离 \ce{Fe^3+} 和 \ce{Al^3+}？
  \item 用两种方法证明 \ce{Fe^3+} 具有氧化性。
  \item 怎样用化学方法除去铜粉中混有的铁粉？
  \item 现有 \ce{CuSO4}、\ce{BaCl2}、\ce{NaOH}、\ce{NaCl} 四种固体，如果你手边除蒸馏水外没有别的试剂，你能用什么方法把它们一一鉴别出来？
  \item 试选用合适的试剂分别将下列难溶物质溶解。
  \begin{tasks}[label=(\arabic*),label-width=16pt]
    \task \ce{Cu(OH)2}
    \task \ce{AgCl}
  \end{tasks}
  \item 利用 \ce{CuSO4}、\ce{BaCl2}、\ce{NaOH}、氨水四种溶液设计一个实验，证明络合物 \ce{[Cu(NH3)4]SO4} 中的 \ce{Cu^2+} 在内界，\ce{SO4^2-} 在外界。
\end{enumerate}
\begin{Theorem}{问题和讨论}
  \ce{AgCl} 与适量的氨水反应可生成络合物，它的化学式是 \ce{AgCl.2NH3}，你怎样通过实验确定它的内界和外界？
\end{Theorem}

\section{甲烷的制取和性质}
\subsection*{实验目的}
\begin{enumerate}[1.]
  \item 学会实验室制取甲烷的方法。
  \item 加深对甲烷性质的认识。
\end{enumerate}
\subsection*{实验用品}
试管、酒精灯、铁架台、水槽、带导管的塞子、 玻璃棒、小烧杯、药匙。

无水醋酸钠、碱石灰、粒状 \ce{NaOH} 固体、石灰水、稀硫酸、$0.01M$ \ce{KMnO4} 溶液。
\subsection*{实验步骤}
\begin{enumerate}[1.]
  \item 甲烷的制取
  \begin{enumerate}[(1)]
    \item 按照\cref{fig:a1} 把仪器安装好，并检查这个装置是否漏气。
    \begin{figure}
      \includegraphics{A-1.pdf}
      \caption{甲烷的制取}\label{fig:a1}
    \end{figure}
    \item 取 \qty{2.5}{g} 无水醋酸钠和 \qty{1.5}{g} 碱石灰 \footnote{在上述混和物中加约 \qty{1}{g} \ce{NaOH} 可提高甲烷的产率。实验室中亦可用 \qty{2.5}{g} 无水醋酸钠，\qty{1.5}{g} \ce{NaOH} 和 \qty{1.5}{g} \ce{CaO} 的混和物来制取甲烷。}，在纸上用玻璃棒混和均匀。把这混和物装入上述装置的试管里。
    \item 小心地加热盛有混和物的试管，等到试管里的空气排出后，用排水取气法收集一试管甲烷。写出制取甲烷反应的化学方程式。
  \end{enumerate}
  \item 甲烷的性质
  \begin{enumerate}[(1)]
    \item 观察甲烷的颜色、状态和闻它的气味。
    \item 检查甲烷纯度后，在导管口点燃甲烷。在甲烷火焰的上方倒放一个干燥的烧杯。注意观察现象。再换一个内壁用石灰水润湿了的烧杯，罩在甲烷火焰上，观察并解释所观察到的现象，写出有关反应的化学方程式。
    \item 把导管下端的弯曲玻璃管取下，换上一根直的玻璃管。把导管插入一个盛有 \qty{2}{ml} $0.01M$ \ce{KMnO4} 溶液（溶液里先滴入少量硫酸）的试管里。通入甲烷后，观察溶液颜色有没有变化。
  \end{enumerate}
\end{enumerate}

\begin{Theorem}{问题和讨论}
\begin{enumerate}[1.]
  \item 盛有反应物的试管的管口为什么要向下倾斜？
  \item 为什么要用无水的醋酸钠？用存放时间较久的碱石灰作实验为什么效果往往不好？
\end{enumerate}
\end{Theorem}

\section{乙烯、乙炔的制取和性质}
\subsection*{实验目的}
\begin{enumerate}[1.]
  \item 学会乙烯、乙炔的实验室制法。
  \item 加深对乙烯、乙炔性质的认识。
\end{enumerate}
\subsection*{实验用品}
试管、酒精灯、铁架台、大试管、分液漏斗、带导管的塞子。

碳化钙、棉花、酒精跟硫酸的混和液\footnote{混和液是由 1 体积酒精 （浓度不得低于 95\%）和 3 体积浓硫酸混和而成。教师应在实验课前配好，供学生实验时使用。}、碎瓷片、溴水、酸性 \ce{KMnO4} 溶液。
\subsection*{实验步骤}
\begin{enumerate}[1.]
  \item 乙烯的制取和性质
  \begin{enumerate}[(1)]
    \item 按照\cref{fig:a2} 把仪器连接好，并检查这一装置是否漏气。
    \item 向试管里倒入酒精跟硫酸的混和液 \qtyrange{6}{8}{ml}，再在这个试管里放少量碎瓷片，以免液态混和物在受热时剧烈跳动。然后，用带导管的塞子塞住试管。用酒精灯加热，使混和物温度迅速上升，这时有乙烯产生。
    \item 使乙烯通入盛溴水的试管，观察有什么现象发生。写出反应的化学方程式。
    \item 使乙烯通入酸性 \ce{KMnO4} 溶液,观察有什么现象发生。
    \item 在导气管管口点燃乙烯，观察乙烯燃烧时的火焰。写出反应的化学方程式。
  \end{enumerate}
  \begin{figure}
    \begin{minipage}[b]{0.48\linewidth}\centering
      \includegraphics{A-2.pdf}
      \caption{制取乙烯的装置}\label{fig:a2}
    \end{minipage}
    \begin{minipage}[b]{0.48\linewidth}\centering
      \includegraphics{A-3.pdf}
      \caption{制取乙炔的装置}\label{fig:a3}
    \end{minipage}
  \end{figure}
  \item 乙炔的制取和性质
  \begin{enumerate}[(1)]
    \item 按照\cref{fig:a3} 把仪器安装好。
    \item 在试管里放入几小块碳化钙。分液漏斗里加入约 \qty{15}{ml} 水。
    \item 向试管里滴水。控制好分液漏斗的滴液速度，使反应不要太快。让生成的乙炔先通入盛有酸性 \ce{KMnO4} 溶液的试管，观察发生的现象。
    \item 使乙炔通入盛有溴水的试管，观察发生的现象。
    \item 点燃生成的乙炔，注意观察乙炔的火焰。写出反应的化学方程式。
  \end{enumerate}
\end{enumerate}

\begin{Theorem}{问题和讨论}
\begin{enumerate}[1.]
  \item 试比较乙烯和甲烷在通入酸性 \ce{KMnO4} 溶液后发生的现象，并从结构上给予解释。
  \item 在制取乙烯的反应里使用浓 硫酸的量很大，考虑一下，在这里除了利用硫酸作催化剂外，还利用了它的什么性质来促进反应的进行。
\end{enumerate}
\end{Theorem}


\section{苯和甲苯的性质}
\subsection*{实验目的}
加深对苯和甲苯性质的认识，学习硝化反应的实验操作。

\subsection*{实验用品}
试管、滴管、酒精灯、铁架台、大烧杯、温度计。

苯、甲苯、植物油、溴水、酸性 \ce{KMnO4} 溶液、浓硫酸、浓硝酸（密度 \qty{1.4}{g/cm^3}）。

\subsection*{实验步骤}
\begin{enumerate}[1.]
  \item 在一个试管里加入 \qty{1}{ml} 苯，在另一个试管里加入 \qty{1}{ml} 水。两个试管里各加几滴植物油，振荡。观察油脂在苯里和水里的溶解性。
  \item 在两个试管里各加入 \qty{1}{ml} 苯。向其中的一个试管里加溴水,向另一个试管里加酸性 \ce{KMnO4} 溶液。振荡后,观察有什么现象发生。
  \item 用跟实验 2 所述相同的方法来试验甲苯对酸性 \ce{KMnO4} 溶液的作用。观察到有什么现象发生？
  \item 在一大试管里加入 \qty{1.5}{ml} 浓硝酸和 \qty{2}{ml} 浓硫酸，待混和酸冷却后，逐渐滴入 \qty{1}{ml} 苯，同时振荡试管，使它混和均匀。然后放在 \qty{60}{\celsius} 的水浴（\cref{fig:a4}）中加热，\qty{10}{min} 后，把试管里的物质倒入盛着大量水的烧杯中去。硝酸、硫酸就溶解在水里，生成的硝基苯成一种黄色的液体聚集在杯底。同时，可以闻到一股苦杏仁的气味。
  \begin{figure}
    \includegraphics{A-4.pdf}
    \caption{苯的硝化反应}\label{fig:a4}
  \end{figure}
\end{enumerate}

实验完毕后，把得到的硝基苯，倒在教师指定的容器中。

\begin{Theorem}{问题和讨论}
\begin{enumerate}[1.]
  \item 在苯的硝化反应的实验里，为什么用水浴而不用酒精灯直接加热？在反应里浓硫酸起什么作用？为什么要等混和酸冷却后再逐渐滴入苯？
  \item 怎样鉴别苯、乙苯（\ce{C6H5C2H5}）和苯乙烯（\ce{C6H5CH=CH2}）？
\end{enumerate}
\end{Theorem}

\section{乙醇和苯酚的性质}
\subsection*{实验目的}
加深对乙醇和苯酚的重要性质的认识；熟练掌握连接简单仪器和检查仪器气密性的操作。

\subsection*{实验用品}
试管、试管夹、酒精灯、药匙、单孔橡皮塞、导管、U 形管、烧杯、铁架台。

乙醇、浓硫酸、铜丝、溴化钠、5\% \ce{NaOH} 溶液、苯酚、稀盐酸、浓溴水、\ce{FeCl3} 溶液。

\subsection*{实验步骤}
\begin{enumerate}[1.]
  \item 乙醇氧化生成乙醛

  在试管里加 \qty{1}{ml }乙醇。把一端弯成螺旋状的铜丝放在酒精灯火焰上加热，使铜丝表面有一薄层黑色的氧化铜生成，立即把它插入盛有乙醇的试管里，如此反复操作几次，注意闻生成的乙醛的气味，并注意观察铜丝表面的变化。写出反应的化学方程式。
  \item 乙醇跟氢溴酸反应
  
  在一个大试管里加入 \qty{2}{ml} 乙醇和 \qty{2}{ml}水。在不断摇动和冷却下慢慢加入 \qty{4}{ml} 浓硫酸，等试管内的混和液冷却到室温后，加入 \qty{3}{g} 研细的 \ce{NaBr}。然后按\cref{fig:a5} 把装置连接好。

用小火隔着石棉网加热，\qty{10}{min} 左右，U 形管中开始出现不溶于水的溴乙烷。
  \begin{figure}
    \includegraphics{A-5.pdf}
    \caption{乙醇跟氢溴酸反应的装置}\label{fig:a5}
  \end{figure}
  \item 苯酚的性质
  \begin{enumerate}[(2)]
    \item 在试管里放入少量苯酚的晶粒，再加一些水，振荡试管里的物质，结果水变浑浊。随后加热苯酚和水的混和物，液体逐渐变得透明。再让液体冷却，观察所起的变化，并解释这些现象。
    \item 向苯酚和水的混和物里注入 5\% \ce{NaOH} 溶液，振荡，观察有什么现象发生？解释这个现象，并写出反应的化学方程式。
    \item 在上述溶液中注入少量的盐酸，解释溶液再度变成浑浊的原因，写出反应的化学方程式。
    \item 向透明苯酚溶液里加入浓的溴水，加到有白色浑浊现象出现为止。解释发生这现象的原因。并写出反应的化学方程式。
    \item 在试管里滴入几滴 2\% 苯酚的稀溶液，把 \ce{FeCl3} 溶液逐滴滴入。观察有什么现象发生？
  \end{enumerate}
\end{enumerate}


\begin{Theorem}{问题和讨论}
\begin{enumerate}[1.]
  \item 实验步骤 1 中，铜丝的加热以及插入乙醇里为什么要反复操作几次？
  \item 现有苯酚的苯溶液，怎样从其中提取苯酚？
\end{enumerate}
\end{Theorem}

\section{乙醛的性质}
\subsection*{实验目的}
加深对乙醛的重要化学性质的认识；学习水浴加热的方法。

\subsection*{实验用品}
试管、试管夹、酒精灯、烧杯。

10\% \ce{NaOH} 溶液、2\% \ce{AgNO3} 溶液、2\% 氨水、乙醛稀溶液、 2\% \ce{CuSO4} 溶液。
\subsection*{实验步骤}
\begin{enumerate}[1.]
  \item 在试管里倒入少量的 \ce{NaOH} 溶液，加热煮沸，把碱溶液倒去后，再用清水冲洗。准备做下述实验。
  \item 银镜反应\quad 在上面洗净的试管里注入 \qty{1}{ml} 2\% \ce{AgNO3} 溶液，逐滴加入 2\% 氨水，边滴边振荡，直到最初生成的沉淀刚好溶解为止。然后，沿试管壁滴入 3 滴稀乙醛溶液，把试管放在盛有热水的烧杯里，静置几分钟后，观察试管内壁有什么现象发生。解释这个现象，并写出反应的化学方程式。
  \item 乙醛被氢氧化铜氧化\quad 在试管里注入 10\% \ce{NaOH} 溶液 \qty{2}{ml}，滴入 2\% \ce{CuSO4} 溶液 \numrange{4}{5} 滴，振荡。然后加入 \qty{0.5}{ml} 乙醛溶液，加热到沸腾，观察有什么现象产生。解释这个现象，并写出反应的化学方程式。
\end{enumerate}

\begin{Theorem}{问题和讨论}
做银镜反应实验之前，为什么要用热的 \ce{NaOH} 溶液洗涤试管？
\end{Theorem}


\section{乙酸乙酯的制取}
\subsection*{实验目的}
\begin{enumerate}[1.]
  \item 加深对乙醇、乙酸的性质及酯化反应的认识；
  \item 练习连接简单仪器。
\end{enumerate}
\subsection*{实验用品}
试管、试管夹、铁架台、酒精灯、单孔橡皮塞、导管、pH 试纸、烧杯。

固体乙酸钠、浓硫酸、乙醇、乙酸、饱和碳酸钠溶液。

\subsection*{实验步骤}
\begin{enumerate}[1.]
  \item 乙酸的制取和性质
  
  在试管里放入 \qtyrange{1}{2}{g} 固体乙酸钠，注入 \qtyrange{1}{2}{ml} 浓硫酸，把连着导管的塞子塞好。把试管固定在铁架台上，用酒精灯加热，试管里有蒸气生成，让蒸气通过导管冷凝成液体并收集于另一试管里（也可以把收集液体的试管放在盛冷水的烧杯中）。闻闻生成的液体有什么气味并用 pH 试纸检验，测定它的 pH 值。
  \item 酯化反应——乙酸乙酯的制取
  
  在一个试管里注入乙醇、乙酸各 \qty{2}{ml}，再慢慢加入 \qty{0.5}{ml} 浓硫酸，按\cref{fig:a6} 把装置连接好。

用小火加热试管里的混和物。把产生的蒸气经导管通到 \qty{3}{ml} 饱和 \ce{Na2CO3} 溶液的上方 \qtyrange{2}{3}{mm} 处，注意观察有透明的油状液体浮在液面上，取下盛有碳酸钠溶液的试管，并停止加热。振荡盛有碳酸钠溶液的试管后静置，待溶液分层后，观察上层的油状液体乙酸乙酯，并可闻到果香气味。写出反应的化学方程式。
\begin{figure}
  
\includegraphics{A-6.pdf}
  \caption{乙酸乙酯的制备}\label{fig:a6}
\end{figure}
\end{enumerate}

\begin{Theorem}{问题和讨论}
\begin{enumerate}[1.]
  \item 用乙酸钠跟浓盐酸能不能制取乙酸？
  \item 在制备乙酸乙酯的实验中，浓硫酸和碳酸钠溶液各起什么作用？
\end{enumerate}
\end{Theorem}

\section{葡萄糖、蔗糖、淀粉和纤维素的性质}
\subsection*{实验目的}
加深对葡萄糖、蔗糖、淀粉和纤维素的重要性质的认识。

\subsection*{实验用品}
试管、酒精灯、烧杯、玻璃棒、试管夹、滴管。

10\% 葡萄糖溶液、2\% 蔗糖溶液、棉花、银氨溶液、\ce{CuSO4} 溶液、\ce{NaOH} 溶液、稀硫酸、浓硫酸、淀粉。
\subsection*{实验步骤}
\begin{enumerate}[1.]
  \item 葡萄糖的还原反应
  \begin{enumerate}[(1)]
    \item 在一个洁净的试管里加入 \qty{3}{ml} 银氨溶液，再加入 \qtyrange{1}{2}{ml} 的 10\% 葡萄糖溶液，充分混和后,在酒精灯上或热水浴中加热。观察发生的现象。
    \item 在另一个试管里，加入 \qtyrange{1}{2}{ml} \ce{NaOH} 溶液，再加 \numrange{3}{4} 滴 \ce{CuSO4} 溶液，可以观察到有 \ce{Cu(OH)2} 沉淀生成。再加入 \qty{2}{ml} 10\% 葡萄糖溶液。加热后，可以观察到有红色的氧化亚铜（\ce{Cu2O}）生成。
  \end{enumerate}

  从上述实验，证明葡萄糖分子含有什么官能团。
  \item 蔗糖的水解反应
  \begin{enumerate}[(1)]
    \item 用上面所说的同样的方法，在试管里制备好 \(\mathrm{{Cu}}{\left( \mathrm{{OH}}\right) }_{2}\) 沉淀,然后加入蔗糖溶液,加热。观察有没有生成氧化亚铜沉淀。
    \item 在试管里加入少量蔗糖溶液，再加入 3-5毫升稀硫酸，然后把这混和物煮沸几分钟，使蔗糖发生水解反应。然后加入 \(\mathrm{{NaOH}}\) 溶液来中和其中的 \({\mathrm{H}}_{2}{\mathrm{{SO}}}_{4}\) 。
    \item 在另一个试管里制备 \(\mathrm{{Cu}}{\left( \mathrm{{OH}}\right) }_{2}\) 沉淀,然后加入已经水解的蔗糖溶液，加热，观察有什么现象发生。
  \end{enumerate}

  写出上述实验各步反应的化学方程式。
  \item 淀粉的水解
  \begin{enumerate}[(1)]
    \item 在试管里加入一些研得很细的淀粉，加入少量水，用力振荡。再加入热水，并煮沸，即制得淀粉溶液。
    \item 在一个试管里加入 \qtyrange{1}{2}{ml} 淀粉溶液，再加入少量新制的 \ce{Cu(OH)2} 沉淀。加热后，观察有没有红色氧化亚铜生成。
    \item 在另一个试管里加入 \qtyrange{1}{2}{ml} 淀粉溶液，再加入 \numrange{3}{5} 滴稀硫酸，煮沸 \qty{5}{min}。加碱中和硫酸后，再加入新制的 \ce{Cu(OH)2} 沉淀。加热后，观察有没有红色氧化亚铜生成。
  \end{enumerate}
  \item 纤维素的水解
  \begin{enumerate}[(1)]
    \item 把一小团脱脂棉或一小块滤纸放入试管里，小心地滴入几滴浓硫酸。用玻璃棒搅动使它变成糊状。再加入 \qty{2}{ml} 水，把试管放在沸水里边加热、边摇动 \qty{5}{min}，直到溶液呈亮棕色为止。最后加碱中和。
    \item 把新制的 \ce{Cu(OH)2} 沉淀，加入上述亮棕色液体里，然后加热到沸腾。观察有没有红色氧化亚铜生成。
  \end{enumerate}

  解释这些现象，并写出反应的化学方程式。
\end{enumerate}

\begin{Theorem}{问题和讨论}
\begin{enumerate}[1.]
  \item 实验步骤 1(1) 里，为什么要使用洁净的试管，怎样使试管洁净？
  \item 通过淀粉和纤维素的水解产物跟新制的 \ce{Cu(OH)2} 的反应，说明淀粉和纤维素在组成上有什么相似的地方？
\end{enumerate}
\end{Theorem}


\section{蛋白质的性质}
\subsection*{实验目的}
巩固蛋白质性质的知识。

\subsection*{实验用品}
试管、酒精灯、试管夹、玻璃棒、滴管。

棉布和毛料各一小条、鸡蛋白、硫酸铵饱和溶液、浓硝酸、10\% \ce{CuSO4} 溶液、蒸馏水、40\% 甲醛溶液、95\% 乙醇溶液。

\subsection*{实验步骤}
\begin{enumerate}[1.]
  \item 蛋白质的灼烧\quad 取棉布上的棉线和毛料上的毛线分别在火焰上灼烧，可以闻到不同的气味。
  \item 蛋白质的加热凝固\quad 在试管里加入鸡蛋白的水溶液，加热。可以观察到蛋白质凝成絮状。把絮状蛋白质取出一些放在水里，观察是否溶解。
  \item 硫酸铵对蛋白质溶解性的影响\quad 在试管里加入 \qtyrange{1}{2}{ml} 蛋白质溶液，然后加入饱和硫酸铵溶液。可以观察到蛋白质沉淀出来。把少量有沉淀的液体倒入另一个盛蒸馏水的试管里，可以观察到沉淀的蛋白质重新溶解。
  \item 重金属盐对蛋白质的凝结作用\quad 在试管里加入 \qty{3}{ml} 蛋白质溶液，随即加入 \qty{3}{ml} 10\% \ce{CuSO4} 溶液。可以观察到蛋白质凝结成沉淀析出。把少量沉淀放入盛蒸馏水的试管里，沉淀不再溶解。
  \item 甲醛和乙醇对蛋白质的凝结作用\quad 在两个试管里各加入 \qty{3}{ml} 蛋白质溶液，随即分别加入 \qty{2}{ml} 甲醛溶液和 \qty{2}{ml} 乙醇溶液。可以观察到蛋白质凝结成沉淀析出。把少量沉淀放入盛蒸馏水的试管里，观察沉淀是否溶解。
  \item 蛋白质的颜色反应\quad 向试管里加入少量蛋白质溶液，然后加入几滴浓硝酸，加微热，即有黄色沉淀析出。这就是黄蛋白反应。
\end{enumerate}

\begin{Theorem}{问题和讨论}
\begin{enumerate}[1.]
  \item 甲醛和乙醇对蛋白质的凝结作用在实际上有什么应用？
  \item 试解释波尔多液为什么能杀灭作物的病菌？
\end{enumerate}
\end{Theorem}


\section{酚醛树脂的制取}
\subsection*{实验目的}
加深对缩聚反应和酚醛树脂制法的认识，巩固苯酚、甲醛性质的知识。

\subsection*{实验用品}
大试管、带玻璃导管的橡皮塞、铁架台、酒精灯、蒸发皿、烧杯、量筒。

苯酚、40\% 甲醛溶液、浓盐酸、浓氨水、酒精。

\subsection*{实验步骤}
\begin{enumerate}[1.]
  \item 在大试管里加入 \qty{2.5}{g} 苯酚，\qty{2.5}{ml} 40\% 甲醛溶液，\qty{1}{ml} 浓盐酸，用带玻璃导管的塞子塞好。在沸水浴中加热。过一会儿就可以看到混和物开始剧烈沸腾。从水浴中取出试管，等反应不再剧烈进行时，继续加热（\cref{fig:a7}），直到混和物生成不溶于水的树脂。从水浴中取出试管，使它冷却，倒去试管里上层的液体，下层物质就是酚醛树脂。观察生成树脂的颜色、状态。
  \begin{figure}
    \includegraphics{A-7.pdf}
    \caption{酚醛树脂的制取}\label{fig:a7}
  \end{figure}
  \item 把盛 \qty{2.5}{g} 苯酚和 \qtyrange{3}{4}{ml} 甲醛的试管放在水浴里加热，加入 \qty{1}{ml} 浓氨水以代替浓盐酸。照实验步骤 1 的操作进行。观察生成树脂的颜色、状态。
\end{enumerate}

实验完毕应及时清洗试管。
如果不容易洗净，可以加入少量乙醇，浸泡几分钟，然后清洗。


\section{有机物熔点、沸点的测定}
\subsection*{实验目的}
学习粗略测定有机物熔点、沸点的方法。

\subsection*{实验用品}
试管、温度计、烧杯、橡皮塞、沸石、三脚架、酒精灯、石棉网、试管夹、自制搅棒。

无水乙醇或丙酮、萘。
\subsection*{实验步骤}
\begin{enumerate}[1.]
  \item 测乙醇（或丙酮）的沸点。
  \begin{enumerate}[(1)]
    \item 在试管里加入 \qtyrange{4}{5}{ml} 无水乙醇，放几粒沸石，插上带塞子（塞子上有豁口一，见\cref{fig:a8}）的温度计，温度计距乙醇液面 \qtyrange{2}{3}{cm}，如图所示。
    \begin{figure}
      \begin{minipage}[b]{0.48\linewidth}\centering
        \includegraphics{A-8.pdf}
        \caption{测定沸点的简易装置}\label{fig:a8}
      \end{minipage}
      \begin{minipage}[b]{0.48\linewidth}\centering
        \includegraphics{A-9.pdf}
        \caption{测定熔点的简易装置}\label{fig:a9}
      \end{minipage}
    \end{figure}
    \item 把试管浸入沸水中\footnote{如果测丙酮的沸点，水浴温度应在 \qty{80}{\celsius} 左右。}，使试管盛乙醇的部分刚好浸在沸水里。注意观察试管中的乙醇和温度计的变化。记下乙醇沸腾时的温度（沸腾时液体和蒸气的温度相同）。沸腾时，温度保持相对稳定。观察片刻后取出试管。等温度下降 \qty{10}{\celsius} 左右，再加热，重复上述实验操作三次，每次记下的温度之差不得超过 \qty{1}{\celsius}。算出三次温度的平均值。这就是测得的乙醇的沸点。
  \end{enumerate}
  \item 测萘的熔点
  \begin{enumerate}[(1)]
    \item 取 \qtyrange{3}{4}{g} 萘放在试管里。
    \item 把事先弯制好的搅棒和温度计一起放入试管里，如\cref{fig:a9} 所示。
    \item 把试管底部浸在沸水中，使萘受热熔化。
    \item 从沸水中取出试管。一面上下移动搅棒，一面仔细观察试管里的萘。当有晶体出现时，记下这时的温度。继续搅动萘，到萘全部凝固为止，注意观察萘全部凝固前、后的温度变化。
    \item 再把试管浸入沸水，如前反复操作三次，记录的每次温度之差不能大于 \qty{0.5}{\celsius}。算出三次温度的平均值，这就是测得的萘的熔点。
  \end{enumerate}
\end{enumerate}

\begin{Theorem}{问题和讨论}
\begin{enumerate}[1.]
  \item 在测乙醇沸点时，为什么不把整个试管浸在水里？为什么塞子上要有个豁口？
  \item 测萘的熔点时，为什么要不停地搅动？
\end{enumerate}
\end{Theorem}

\section{实验习题}
\subsection*{实验目的}
巩固一些重要有机物性质的知识。

\subsection*{实验习题}
\begin{enumerate}[1.]
  \item 领取一种溶液，试用两种方法证明它不是醇而是苯酚的溶液。
  \item 三个试管里，分别盛有乙醇、乙醛和乙酸，鉴别哪个试管里是哪种物质。
  \item 怎样鉴别苯、乙醇和液态不饱和烃？
  \item 有两种白色粉末，它们是氯化钠或醋酸钠，怎样区别它们？
  \item 怎样用实验方法证明硬脂酸有酸性？
  \item 怎样检验葡萄、熟苹果中都含有葡萄糖？
  \item 下列几组化合物，怎样用实验方法鉴别？
  \begin{tasks}[item-indent=5em](2)
    \task 甲烷跟乙烯，
    \task 苯跟甲苯，
    \task 苯跟苯酚溶液，
    \task 乙酸跟苯酚溶液，
    \task* 葡萄糖溶液跟蔗糖溶液。
  \end{tasks}
  \item 试用实验方法证明甲酸分子里既有羧基又有醛基。
  \item 怎样鉴定马铃薯里含有淀粉？
  \item 燃烧聚乙烯、聚氯乙烯、聚苯乙烯各一小块，各有什么现象？根据各高分子的组成进行分析。
\end{enumerate}

\section*{选做实验\texorpdfstring{\quad}{ }纸上层析}
\addcontentsline{toc}{section}{选做实验\texorpdfstring{\quad}{ }纸上层析}
\subsection*{实验目的}
初步了解纸上层析的方法

\subsection*{实验用品}
培养皿、表面皿（盖培养皿用）、滴管、圆形滤纸。

\ce{FeCl3} 和 \ce{CuSO4} 的混和液\footnote{\ce{FeCl3} 和 \ce{CuSO4} 的混和液。饱和 \ce{FeCl3} 溶液$:$饱和 \ce{CuSO4} 溶液 ${}=1:1$（体积比）。}、丙酮和盐酸的混和液\footnote{丙酮和盐酸的混和液。丙酮和 $6N$ 盐酸的体积比是 $9:1$。}、浓氨水、甲基橙和酚酞的混和酒精溶液\footnote{甲基橙和酚酞的混和酒精溶液。把 \qty{0.1}{g} 甲基橙 和 \qty{0.1}{g} 酚酞溶在 \qty{10}{ml} 60\% 的酒精溶液里。}、正丁醇和氨水的混和液\footnote{正丁醇—氨水混和液。取 \qty{10}{ml} 正丁醇（如没有正丁醇 可用乙醇代替）和 \qty{4}{ml} 浓氨水充分混和，待分层后，取用上层液体。}、饱和 \ce{Na2CO4} 溶液。
\subsection*{实验步骤}
\begin{enumerate}[1.]
  \item \ce{Fe^3+} 和 \ce{Cu^2+} 的纸上层析
  \begin{enumerate}[(1)]
    \item 在滤纸中心滴一滴 \ce{FeCl3}--\ce{CuSO4} 混和液，液滴直径约 \qty{0.5}{cm}，放置待干。
    \item 把上述滤纸剪开一小条（如\cref{fig:a10}），剪到点样痕迹，小条宽 \qtyrange{2}{3}{mm}。把小条向上折成直角。
    \begin{figure}
      \begin{minipage}[b]{0.48\linewidth}\centering
        \includegraphics{A-10.pdf}
        \caption{纸上层析用滤纸}\label{fig:a10}
      \end{minipage}
      \begin{minipage}[b]{0.48\linewidth}\centering
        \includegraphics{A-11.pdf}
        \caption{纸上层析的展开操作}\label{fig:a11}
      \end{minipage}
    \end{figure}
    \item 把折好的滤纸放在一盛有 \qty{10}{ml} 丙酮—盐酸混和液的培养皿上，使剪开的小条浸入混和液（如\cref{fig:a11}）。
    \item 待滤纸中心的点样痕迹的直径展开到滤纸直径的三分之二左右，取下滤纸，放在一盛有浓氨水的培养皿上，用氨气熏。注意不要使氨水直接浸到滤纸上。
    
    一会儿，可以看到滤纸上出现两个有色的同心圆环，外圈是黄色环，内圈是蓝色环。
  \end{enumerate}
  \item 甲基橙和酚酞的纸上层析
  \begin{enumerate}[(1)]
    \item 在滤纸中心滴一滴甲基橙—酚酞的混和液（或围绕滤纸中心，距中心约一厘米处滴几个点），液滴直径 \qty{0.5}{cm}，放置待干。
    \item 在上述滤纸中心扎一小洞，插上一根滤纸捻成的细纸芯。
    \item 把插好滤纸芯的滤纸放在盛有正丁醇—氨水混和液的培养皿上，使纸芯浸入混和液。随着液体的渗吸，滤纸中心的黄色点样痕迹向外逐渐扩大，成为黄环。
    \item 待黄环半径扩展到滤纸半径的二分之一时，取下滤纸，拔下纸芯。等滤纸稍干后，喷上饱和 \ce{Na2CO3} 溶液，于是在甲基橙黄环外围出现一个酚酞的红环。
  \end{enumerate}
\end{enumerate}

